\documentclass{article}
\usepackage{graphicx} % Required for inserting images

\title{TP1}
\author{Edgar Israel Vargas Arce}
\date{March 2026}

\begin{document}

\maketitle

\section{Introduction}

\section{Demostración}

Se desea determinar una expresión para el volumen de líquido contenido en un tanque esférico de radio $R$, en función de la altura $h$ medida desde la base del tanque. Para ello, se emplea el método de integración utilizando secciones transversales (discos).

Se considera una esfera de radio $R$, centrada en el origen de un sistema de coordenadas. Por lo tanto, la esfera se extiende desde $y = -R$ hasta $y = R$. Su ecuación es:

\begin{equation}
x^2 + y^2 = R^2
\end{equation}

Despejando $x$, se obtiene el radio de cada sección transversal en función de $y$:

\begin{equation}
x = \sqrt{R^2 - y^2}
\end{equation}

Cada corte horizontal de la esfera genera un disco de radio $x$ y, por tanto, de área:

\begin{equation}
A(y) = \pi x^2 = \pi (R^2 - y^2)
\end{equation}

Para calcular el volumen del líquido, se integran estas áreas desde la base del tanque hasta la altura del líquido.

Dado que el origen está en el centro de la esfera:
\begin{itemize}
\item La base del tanque corresponde a $y = -R$
\item La superficie del líquido corresponde a $y = h - R$
\end{itemize}

Esto se debe a que la altura $h$ se mide desde la base, por lo que al trasladarla al sistema centrado en el origen, se debe restar $R$.

Por lo tanto, los límites de integración son:

\begin{equation}
V = \int_{-R}^{h-R} \pi (R^2 - y^2)\,dy
\end{equation}

Resolviendo la integral:

\begin{equation}
V = \pi \int_{-R}^{h-R} (R^2 - y^2)\,dy
\end{equation}

\begin{equation}
\int (R^2 - y^2)\,dy = R^2 y - \frac{y^3}{3}
\end{equation}

Evaluando en los límites:

\begin{equation}
V = \pi \left[ R^2 y - \frac{y^3}{3} \right]_{-R}^{h-R}
\end{equation}

Finalmente, al simplificar la expresión anterior, se obtiene:

\begin{equation}
V = \frac{\pi h^2 (3R - h)}{3}
\end{equation}



\end{document}
